\section*{Complementi di Analisi Matematica. Foglio n.2 (19/9/2025)}
\addcontentsline{toc}{section}{Complementi di Analisi Matematica. Foglio n.2 (19/9/2025)}

\textbf{Esercizi su topologia e funzioni continue}

\begin{enumerate}
	\item Sia $f: A \to \mathbb{R}^m$, ove $A \subset \mathbb{R}^n$ e sia $F(x) = (x, f(x))$ per ogni $x \in A$, $F: A \to \mathbb{R}^n \times \mathbb{R}^m$ la funzione grafico di $f$. Provare che $f: A \to \mathbb{R}^m$ è continua se e solo se $F: A \to \mathbb{R}^n \times \mathbb{R}^m$ è continua.
	
	\item Data la funzione $f(x,y) = \frac{1}{x^8 - y^8}$,
	\begin{enumerate}
		\item determinare il più grande insieme $A \subset \mathbb{R}^2$ dove $f$ è definita,
		\item determinare il più grande insieme $A \subset \mathbb{R}^2$ dove $f$ è definita e continua.
	\end{enumerate}
	
	\item Tracciare il grafico della funzione $f: \mathbb{R}^2 \to \mathbb{R}$, $f(x,y) = \sqrt[4]{x^2 + y^2}$ e stabilire se tale insieme è chiuso, aperto, limitato, compatto o connesso per archi.
	\textit{Sugg.} Si consideri la funzione grafico $F(x,y) = \left(x, y, \sqrt[4]{x^2 + y^2}\right)$, che è continua.
	
	\item Si tracci l'insieme $A = \{(x,y,z) \in \mathbb{R}^3 : x^2 + 2y^2 + z \leq 8\}$, stabilendo se è connesso per archi.
	\textit{Sugg.} Si consiglia di tracciare prima il grafico di $f(x,y) = 8 - x^2 - 2y^2$.
	
	\item Stabilire se $f(x,y) = \frac{\sin(xy)}{\sqrt[3]{x^2 + y^2}}$ è continua su $A = \mathbb{R}^2 \setminus \{(0,0)\}$.
	
	\item Consideriamo $f: S \to \mathbb{R}$, $f(x,y,z) = \frac{x^2 + y^2}{1 + z^2 - x^2 - y^2}$, ove
	\[
	S = \left\{(x,y,z) \in \mathbb{R}^3 : x^2 + y^2 = z^2 \leq 1\right\}.
	\]
	\begin{enumerate}
		\item Tracciare un grafico qualitativo di $S$.
		\item Studiare l'esistenza del massimo e del minimo di $f$ su $S$.
		\item Determinare l'immagine $f(S)$.
	\end{enumerate}
	
	\item Consideriamo $f: B(0,2) \to \mathbb{R}$ definita come $f(x,y) = 1 - x^2 - y^2$.
	\begin{enumerate}
		\item Tracciare il grafico di $f$.
		\item Stabilire se $f$ ha massimo e minimo assoluti o relativi.
		\item Determinare l'immagine di $f$, $\sup_{B(0,2)} f$ e $\inf_{B(0,2)} f$.
	\end{enumerate}
	
	\item Consideriamo $f: [0,1]^2 \to \mathbb{R}$ e $f(x,y) = -\sin^2(x-y)$. Determinare massimo e minimo di $f$ utilizzando solo disuguaglianze.
	
	\item Fissiamo $r > 0$, $B(0,r) \subset \mathbb{R}^n$ ed $f: B(0,r) \to \mathbb{R}$ con $f(x) = 1 - |x|^2$. Determinare l'immagine di $f$, calcolando $\sup_{B(0,r)} f$, $\inf_{B(0,r)} f$ e stabilendo se esistono massimo e minimo di $f$ su $B(0,r)$.
\end{enumerate}